\chapter*{Introduction and scope}
\addcontentsline{toc}{chapter}{Introduction and scope}
\setcounter{section}{0}\renewcommand{\thesection}{I.\arabic{section}}

% --- source: tex/s01_intro#0 ((head))


% --- source: tex/s01_intro#1 (Introduction)
\section{Introduction}
\label{sec:intro}

\subsection{The question}
\label{sec:question}

Can the mathematical structures underlying spacetime, electromagnetism and matter be understood as different manifestations of a common
algebraic foundation, rather than introduced as independent structures from the outset?

The present work investigates this question through the real split-signature Clifford algebra
\[
\Cl(3,3)\cong M_8(\R)
\]
and its graded extensions $\Cl(3n,3n)$. The objective is to determine which structures follow from specified algebraic data, which require
additional assumptions, which are mutually compatible and which are obstructed. Representing a known theory within the algebra is
distinguished from deriving a restriction on its physical content.

\paragraph{The algebraic prototype and its temporal interpretation.} We begin with a six-dimensional real vector space
\[
V=V^+\oplus V^-,\qquad \eta=\operatorname{diag}(+,+,+,-,-,-),
\]
where $V^+$ and $V^-$ each have dimension three. A prescribed pairing identifies each positive basis direction with a corresponding
negative direction. We assign a temporal interpretation to the positive triple and a spatial interpretation to the negative triple as part
of the proposed spacetime reading.

One motivation for this choice is the six-dimensional space of electromagnetic field strengths in four-dimensional spacetime, organised
relative to an observer into three electric and three magnetic components. This motivates a comparison of representations; it does not
establish the existence of three physical time coordinates. The displacement and field-strength readings use distinct Lorentz actions, and
their compatibility must be examined explicitly (\cref{ch:ch03_readings}).

A related construction uses the three commuting pole operators
\[
\omega_i=e_ie_{i+3},\qquad \omega_i^2=\id .
\]
Their joint eigenspaces are eight one-dimensional subspaces of the real spinor module. Following the principal convention of
\cref{r3:sec:main}, the field components are assigned to six of these octants:
\begin{center}
\begin{tabular}{@{}cccc@{}}
\toprule
component & pole eigenvalues & component & antipodal pole eigenvalues\\
\midrule
$E_x$ & $(-,+,+)$ & $B_x$ & $(+,-,-)$\\
$E_y$ & $(+,-,+)$ & $B_y$ & $(-,+,-)$\\
$E_z$ & $(+,+,-)$ & $B_z$ & $(-,-,+)$\\
\bottomrule
\end{tabular}
\end{center}
The electric components occupy the three octants of Fock weight one and the magnetic components their antipodal octants of Fock weight
two. The sign triples label joint eigenspaces of the pole operators, not metric signatures. The assignment is a constructed interpretation
rather than a unique consequence of the algebra, and its electromagnetic transformation law is distinct from the interior Dirac spinor
action on the same underlying real space.

For the observational spacetime construction, selecting a unit clock $\hat n\in V^+$ gives the Lorentzian four-plane
\[
E_{\hat n}=\R\hat n\oplus V^- .
\]
Its orthogonal complement is a positive two-plane. The existence of this splitting is algebraically straightforward; its physical
selection, the treatment of the two surplus directions and causal consistency require further analysis. The signature $(3,3)$ is
therefore adopted as a balanced prototype with three prescribed pairs. Whether additional requirements characterise this choice uniquely
remains a mathematical question.

\paragraph{From a single prototype to constituent multiplicity.} The algebra $\Cl(3,3)$ supplies the single-constituent starting point. With
the specified clock and representation data it supports a Lorentzian sector and a real Dirac module, and these constructions can be defined
before a second constituent is introduced. We investigate whether repeated prototypes provide a useful mathematical setting for composite
systems, multiple-particle descriptions and alternative collective states. The factors are isomorphic copies, and their initial labels are
not assigned an intrinsic physical identity; distinctions may enter through states, relations, selected structures or interactions, and
whether factor permutations are physical symmetries depends on those additional data. The extensions are formed through orthogonal direct
sums and graded tensor products,
\[
\Cl(3n,3n)\cong\widehat{\bigotimes}_{k=1}^{n}\Cl(3,3).
\]
This classical identity supplies a systematic way to enlarge the representation space while retaining copies of the prototype. The integer
$n$ counts algebraic constituents. An $n$-particle interpretation additionally requires physical state spaces, kinematics, dynamics and
exchange properties; the tower does not supply particle number or statistics by definition.

\paragraph{Constructive method and compatibility.} We begin with the algebraic data, construct the geometric and physical readings they
admit, and examine their behaviour under extension. Alternative readings are compared through their representations, invariant
structures, normalisation conditions and possible dynamics, and a construction in one sector transfers to another only when the necessary
map and compatibility conditions are supplied. Known physical theories provide benchmarks: familiar identities and representations are
identified as inherited results, and explicit realisations, conditional classifications and obstructions are stated with their
assumptions. The research objective is to determine whether compatible structures can be assembled into a common formulation with fewer
independent physical inputs.

\paragraph{The role of a common action.} Dynamics is introduced through sector-specific actions whose assumptions and scope can be stated
explicitly. These actions provide candidate models and consistency tests; their separate existence does not establish a common variational
origin. The construction of a universal parent action is deferred while the available representations, kinematic identities and
compatibility constraints are investigated. A nearer objective is a common action for a specified compatible subset of sectors, with
controlled constraints, causal evolution and physical degrees of freedom. Whether such an action can incorporate all the proposed readings
remains open. A satisfactory common formulation would need to address the selection of the observational spacetime sector, the matter and
gauge content, the dynamics of the metric and the stability of the resulting physical states. No universal action is claimed in this
manuscript.

\paragraph{The governing criterion.} Throughout the investigation, the assumptions entering a result are distinguished from its
consequences, and mathematical constructions are distinguished from their physical interpretations. Both successful compatibility results
and demonstrated obstructions are part of the outcome. The central question is whether the common algebraic setting reduces the number of
independent physical assumptions while preserving mathematical consistency. Convergence toward such a formulation is the objective;
identifying precisely where proposed combinations fail is equally part of the research.

\paragraph{Related work.} The classification of real Clifford algebras and their gradings is classical
\cite{LawsonMichelsohn,Lounesto,Porteous,MosnaMirallesVaz}. Several programmes extract Standard-Model structure from division and Clifford
algebras \cite{Furey,Gresnigt,ChamseddineConnes}, and the representation-theoretic limits of such embeddings are discussed in
\cite{DistlerGaribaldi}. Krasnov realises one generation on the forms of a doubled split space \cite{Krasnov18}. \Cref{ch:ch14_related}
compares the constructions of this manuscript with these and other approaches, and the comparisons that concern a single statement are made
where the statement occurs.

% --- source: tex/s01_intro#3 (A deflationary principle)
\subsection{A deflationary principle}
\label{sec:deflation}

One theorem, proved in \cref{sec:dirac}, governs everything else. Take the four interior gamma matrices
$\gamma^\mu=(e_1,e_4,e_5,e_6)$ and the operator $J=e_2e_3$, which commutes with all four and squares to $-\id$. Then the
algebra generated over $\C$, with $J$ playing the role of $i$, by the four gamma matrices is the whole of $M_4(\C)$. The
real Dirac theory built on $\Cl(3,3)$ is therefore ordinary complex Dirac theory written in real variables, exactly and
not in a limit. A change of variables predicts nothing, so a physical claim about the arena carries information only if it
could have come out otherwise, that is, only if it adds a constraint beyond a change of notation or a freely chosen embedding. This is a
requirement on physical claims, not a definition of mathematical content, and we apply it throughout. Facts that are true of Dirac theory in any notation, such as
$g=2$ at tree level, current conservation or the Ward identity, are \emph{inherited}; they are reproduced here only where
they are needed to state something else. The statements we report are of a different kind: which bilinears are admissible,
which generators may be gauged, how large the internal centraliser is at each rung, what the faithful gauge algebra of a
projected block is, which potentials can select a vacuum, and which constructions are impossible under stated
hypotheses.

Every claim in this paper belongs to one of four classes, and we say in the text which one. A result is \emph{fixed by the
algebra} when it is an identity or a theorem proved from explicitly stated hypotheses. It is \emph{constructed} when the algebra admits it
and we exhibit it, without anything selecting it among alternatives. It is \emph{obstructed} when it is impossible under
hypotheses that we state, and those hypotheses are part of the result. It is \emph{open} when it is neither established
nor refuted. Three further labels say where a statement comes from or how far it reaches: \emph{inherited} marks a known result of the
literature, used with its hypotheses; \emph{restated} marks a result proved earlier in this programme and re-verified here; \emph{interpretive}
marks a physical reading of a fixed statement, which the mathematics does not decide. No statement is a prediction of a specified physical
model unless it says so. \Cref{ch:ch13_compatibility} applies the same classes to the links between results, with five statuses:
\emph{established} (fixed, or inherited and used within its hypotheses), \emph{constructed}, \emph{conditional} (holding under a stated
assumption, an adopted action or a chosen datum), \emph{obstructed} and \emph{open}.

Every blade of the realisation used here is a signed permutation matrix, so products, commutators, traces and ranks of blades and of their
integer combinations are integer computations. Statements with irrational coefficients, such as those of the symmetric clock
$(e_1+e_2+e_3)/\sqrt3$, parameter-dependent spectra, minimisations and numerical ranks are not, and the numerical checks are identified as such
(\cref{ch:ch15_methodology}). Most sections have a derivation script and a separately written referee script; two sections were refereed by
reading only, and the photon scripts \texttt{p1} and \texttt{p2} have no referee script. Every script prints one \texttt{PASS} or \texttt{FAIL}
line per check and exits with status zero only if every check passes. How the statements were verified is described in \cref{app:methods}.

% --- source: tex/s01_intro#4 (Summary of results)
\subsection{Summary of results}

\paragraph{The arena (\cref{sec:arena,sec:clock,sec:dirac,sec:discrete}).} A bivector $e_ae_b$ squares to $-\id$ when its
two directions have the same sign and to $+\id$ when they have opposite signs. The first kind generates a compact $U(1)$, which can be
gauged once a connection, a kinetic term and a matter coupling are supplied; the second generates a boost, which is not unitary for a positive
form on the module. As operators on the spinor module the first are complex structures and the second have real $\pm1$ eigenspaces, while
conjugation by either kind is an involution of the algebra. The three \emph{pole operators}
$e_ie_{i+3}$ are commuting involutions grading the algebra by $\Z_2^3$, and they span a maximally split Cartan
subalgebra of $\so{3,3}$. The pairing has two square roots: a para-complex structure $\cK$ with $\cK^2=+\id$, which grades,
and a complex structure $G$ with $G^2=-\id$, which is self-adjoint for $\eta$ and, as we show later, binds (gives bounded classical orbits of) the interior
dynamics and acts as a Hodge star. A clock is a unit vector of the positive triple; the symmetric clock
$\widehat W=(e_1+e_2+e_3)/\sqrt3$ is the unique permutation-invariant point and the entropy maximum of the clock simplex
under the Fisher--Rao metric, once clocks are described by probability weights; that description is a choice, and every unit clock is
related to $\widehat W$ by a rotation. Conversion along any single clock is invisible in the equations of motion, so this does not
single out $\widehat W$.
The imaginary unit is selected rather than adjoined: inside the joint commutant of the interior Lorentz action and the
internal bivectors, the only square roots of $-\id$ that commute with the gamma matrices are $\pm J$. Charge conjugation is
the surplus plane $\langle e_2,e_3\rangle$ and coincides with the Majorana structure; time reversal is forced to square to
$-\id$; the kinetic forms alone fix the internal symmetry algebra $\uu1\oplus\su2$ of the free massless action, which becomes a gauge
algebra only once a connection is introduced and which the mass form cuts to the $\uu1$ of $J$; and for the real representation of the
canonical anticommutation relations and the observables constructed in \cref{ch:ch02_lorentz}, the quantum state space must be formed by
complexifying from outside.

\paragraph{The tower and its matter (\cref{sec:tower,sec:projectors,sec:pstar,sec:charges,sec:actions}).} The bivectors
commuting with the interior Lorentz action of $\Cl(3n,3n)$ form $\so{3n-1,3n-3}$, nested along the tower, of dimensions
$1,28,91,190,\dots$. At one constituent this centraliser is the single bivector $J$, the only compact $U(1)$ candidate of the
interior structure; identifying it with the photon would need a connection, a kinetic term, a matter coupling and physical degrees of
freedom, and the electromagnetic generator is constructed only in \cref{ch:ch08_photon}. At three constituents we construct an orthogonal projector
\[
P_*=\tfrac18\,(\id+\Gamma\chi_W)(\id-J\chi_C)(\id-Jq_A)
\]
of rank $64$ which commutes with the Lorentz action and with $J$, is axial in the sense that $P_*\gamma^\mu P_*=0$, keeps the whole of
$\su4$ acting faithfully, and carries $(\mathbf4,\mathbf2,\mathbf1)_L\oplus(\mathbf4,\mathbf1,\mathbf2)_R$: one Pati--Salam
family with a right-handed neutrino. Its two central abelian factors cannot be gauged on this matter. Charges built from
three commuting involutions are conjugately paired whenever they are anomaly free, so hypercharge cannot be read off from
anomaly cancellation inside that class. When the Lorentz action is instead taken diagonally from all three constituents,
the internal algebra becomes the entirely compact $\so6\oplus\so3$; the copy algebra and the Lorentz algebra form a
reductive dual pair satisfying $C_{\rm cop}+C_\Delta=\tfrac{15}2$, both inherited \cite{Howe89,Neergard19}, and the module is vector-like. A single theorem gives
the candidate gauge algebra of every Lorentz action on the tower, and with it a rung-by-rung census of the vector bosons
the tower can carry.

\paragraph{Gauge fields and the electromagnetic reading (\cref{sec:fields,sec:emfield,sec:spin,sec:orth,sec:energy,sec:massbreaks}).} A varying
projector induces a connection whose available curvature algebra is far too large, so a parent action would have to cut
the projector's freedom rather than build a gauge field. The centraliser of the complex structure $G$ in $\so{3,3}$ is a
Lorentz algebra, under which $\R^{3,3}$ \emph{is} $\Lambda^2\R^{1,3}$: boosts act by $\delta\bE=\bv\times\bB$,
$\delta\bB=-\bv\times\bE$, $G$ is the Hodge star, and the two invariants $|\bB|^2-|\bE|^2$ and $\bE\cdot\bB$ are the real
and imaginary parts of one holomorphic quadratic form. This reading excludes the interior four-plane on which the matter
is built. Covariance fixes the form of Maxwell's equations, because the vector module occurs in $\Lambda^2\otimes V$ with
multiplicity exactly two, and fixes nothing else. The photon's spin one appears in two places, related by the Klein
correspondence $V\cong\Lambda^2(S^+)^*$, and at the second rung both can be coupled to a Dirac block in a single gauge
theory. The orthogonality of $\bE$ and $\bB$ is the reality of the invariant $q=z\cdot z$; its failure breaks parity and time
reversal, is never spontaneous for the Maxwell energy, and does not generate mass. The energy of a field point is one undivided, duality-invariant
quantity $E_6=\tfrac12|z|^2$; its rest part $E_0=\tfrac12|q|$ is the infimum over boosts, attained for non-null fields only, and
$E_6^2=E_0^2+|\mathbf S|^2$ is the field-point form of $E^2=p^2c^2+m^2c^4$, an identity of a field point and not a particle mass; $c$ enters only as
the conversion between the two triples. $\mathrm{Spin}(3,3)$ keeps the Dirac mass form, so
a mass breaks no spacetime symmetry, and particle and antiparticle are the two senses of one rest rotation.

\paragraph{Interior dynamics (\cref{sec:shifted,sec:neutral,sec:closure,sec:restclosure,sec:bfield}).} The Lagrangian of a single shifted operator
$\partial_\tau+\kappa g$ factorises exactly when the shift is symmetric, and then it describes an inverted oscillator; for
every constant shift some solution is unbounded, and with a positive form no solution with nonzero momentum recurs. In the
split form the same factorisation admits bound motion, and in the factorised shift model it does so exactly in neutral signature, with the
complex structure $G$ as the shift. The force law is a dual transport which is a Lie--Poisson equation only on an extension of the
frame; within the closure ansatz that extension is unique, its rotor angle is the duality angle of the momentum, and its free motion is a pendulum
that drifts the frame except on null data. The rest energy sets the amplitude of that pendulum; a constant rate and a bounded frame
coexist only on null data or on the balance surface of several constituents, and identifying the rest rotation of the Dirac field
with the rate of the closure, $mc^2=\hbar\kappa$, is an added hypothesis. The antisymmetric part of a shift is a connection: the part
commuting with the binding root is removed by a moving frame, and no connection binds on its own.

\paragraph{The vacuum (\cref{sec:subfields,sec:kernels,sec:selection,sec:matter,sec:goldstone}).} At one constituent we
classify the vacuum of every coordinate three-plane in four senses. At three constituents the space of Lorentz-invariant
mass kernels of the chiral block has dimension $272$ and, with the complex structure $\widetilde{\mathcal C}$, decomposes as
$(\overline{\mathbf{10}},\mathbf3,\mathbf1)\oplus(\mathbf{10},\mathbf1,\mathbf3)\oplus(\mathbf{15},\mathbf2,\mathbf2)\oplus
2(\mathbf6,\mathbf1,\mathbf1)\oplus(\mathbf1,\mathbf2,\mathbf2)$; the Standard-Model invariant kernels form one complex
line, the right-handed neutrino Majorana kernel. We give the complete renormalizable potential of the
$(\mathbf{10},\mathbf1,\mathbf3)$ component, including a holomorphic quartic that violates baryon number by two units,
and sufficient conditions, together with separate necessary conditions, under which the Standard-Model orbit is its global minimum (the two
do not coincide, and numerics support, but do not locate exactly, the transition between them); within the adopted leading-logarithm
model we compute the radiatively generated spectrum of all $271$ angular modes, the Goldstone content and the discrete symmetries of the
vacuum, and, for a given bidoublet value in the kernel model, an electroweak stage with the Georgi--Jarlskog factor and a seesaw.

\paragraph{Curvature and the photon (\cref{sec:curved,sec:curvop,sec:transport,sec:photon}).} In a curved arena the
pairing becomes an almost para-Hermitian structure. An explicit six-dimensional metric with a normal $U(1)$ has a
closed-form Einstein tensor whose tangential block on the four-dimensional slice is the Einstein tensor of the slice, while
relative to the ambient Einstein--Hilbert action and the stated ansatz the normal field strength gravitates with negative energy, as
the signature forces. In the field-strength reading the
curvature operator of a four-metric lives on the frame, the first Bianchi identity is the single condition
$\Tr(G\cR)=0$, and the duality rotation of the curvature ends at the double dual, which contracts to the Einstein tensor.
These are identities of given metrics; no gravitational action is derived. The projector current sources a massive normal vector, and a massless photon requires a vacuum that commutes with the
charge; such a vacuum exists at the second rung and not at the first. At the third rung the normal charge is half the fermion
number on the chiral block and anomalous, and the electric charge is the photon candidate. In this projector model the selected vacuum is electroweak symmetric and electroweak
breaking is not obtained; the electroweak stage of the kernel model belongs to a different model, and a common realisation of the two is not
established.

\paragraph{Causal structure (\cref{sec:baseline,sec:selector,sec:typesel,sec:constraints}).} For unrestricted evolution, a field propagating
freely in six dimensions of signature $(3,3)$ is ill-posed and has no ground state, and its time-like cone has no past and future;
constrained ultrahyperbolic problems can be well posed \cite{CraigWeinstein09}.
Choosing a Lorentzian four-plane inside $\R^{3,3}$ is the same as choosing its definite orthogonal two-plane, represented by a normalised simple
Clifford bivector whose two signs encode an orientation, and that bivector is a complex structure, and at one constituent the
light cone of that plane is the set of chiral spinors squared with the help of $J$. For an unconstrained adjoint scalar, no bounded-below
invariant polynomial potential of degree below eight isolates that choice, and every invariant bilinear kinetic term has ghost directions. The
surplus-time constraints follow from a gauge symmetry of the worldline that is chosen, not derived, and is shown for a test worldline only.

% --- source: tex/s01_intro#5 (Conventions)
\subsection{Conventions}
\label{sec:conventions}

Spacetime signature is $(+,-,-,-)$, units are Heaviside--Lorentz with $\hbar=c=1$ unless stated. A \emph{constituent} is
a tensor factor $\Cl(3,3)$ of the tower; the word never refers to a particle, and three constituents do not describe three
particles. Anomaly coefficients are chiral traces per Weyl component with every field written left-handed. In the
field-strength reading of $\R^{3,3}$ the vector $X=\sum_i(B_ie_i+E_ie_{i+3})$ is identified with the two-form $F$ of
components $F_{0i}=E_i$, $F_{jk}=\epsilon_{ijk}B_i$; with this orientation $\bB$ is minus the textbook magnetic field,
which we write $\Btext=-\bB$, while $\bE$ is the textbook electric field. Formulas taken from the literature are written in
textbook variables, and where a sign matters the variables are named. Compact Lie algebras are realised by skew operators
and the corresponding charges by the symmetric operators $\hat t=-Jt$.

\subsection{Scope}
\label{sec:fm:scope}

The manuscript is a mathematical study of one algebraic arena and of the constructions it admits. It contains:
the representation theory of $\Cl(3n,3n)$ relevant to Lorentz structures, complex units, centralisers and projected blocks; two readings of
$\R^{3,3}$, as displacements and as field strengths, with their Lagrangians and energies; an interior dynamics and its rigid-body form; the
chiral block, its charges and its mass kernels, with a conditional vacuum selection; the identification of a photon candidate; dynamics
between constituents; classical and observational consistency checks of stated readings; curvature identities; and the causal structure of
the arena. It does not contain a derivation of gravitational dynamics, a derived action for any sector, a spectral-triple formulation,
quantised interacting fields, or predictions of a specified physical model. Where a statement is a reading rather than a theorem, it is
marked interpretive. The colour and Standard-Model representation results are kept; no work on a gluon identification is included.

\subsection{Notation}
\label{sec:fm:notation}

\Cref{tab:fm:notation} lists the symbols used throughout. A symbol in the table has one meaning in the whole manuscript; symbols local to a
chapter are defined there, and those of the matter, vacuum and photon chapters are collected in \cref{tab:app:symbols}. Units of the two
readings are in \cref{tab:app:units}.

\begin{table}[t]
\centering
\small
\begin{tabular}{@{}>{\raggedright\arraybackslash}p{4.4cm}>{\raggedright\arraybackslash}p{8.4cm}>{\raggedright\arraybackslash}p{1.6cm}@{}}
\toprule
symbol & meaning & chapter\\
\midrule
$V=\R^{3,3}$, $\eta$, $V^\pm$ & the arena, $\eta=\operatorname{diag}(1,1,1,\allowbreak-1,-1,-1)$, temporal and spatial triples & 1\\
$e_1,\dots,e_6$; $\Cl(3n,3n)$, $S_n$ & generators; tower algebra on $S_n=(\R^8)^{\otimes n}$, $n$ the number of constituents & 1, 5\\
$\omega_i=e_ie_{i+3}$, $\varpi$ & pole operators, their sum & 1\\
$\cK$, $G$ & grading root ($\cK^2=\id$); binding root ($G^2=-\id$, $\eta$-self-adjoint) & 1, 3\\
$T=e_1e_2e_3$ & temporal volume, the Dirac mass form & 2\\
$J$, $J(\hat n)$; $\Gamma=J\Omega$ & complex structure (imaginary unit); chirality & 2, 6\\
$\hat n$, $\widehat W$, $E_{\hat n}$ & clock, symmetric clock, observable four-plane $\R\hat n\oplus V^-$ & 2\\
$\lint$, $\lG$ & interior Lorentz algebra; Lorentz algebra centralising $G$ & 2, 3\\
$\bE$, $\bB$, $z$, $q$; $E_6$, $E_0$ & fields ($\bB=-\Btext$), $z=\bB+i\bE$, $q=z\cdot z$; undivided and rest energy & 3\\
$\Theta^{\mu\nu}$ & stress--energy tensor & 3\\
$\kappa=c/\ell$, $\tau$ & rate and parameter of the interior dynamics & 4, 12\\
$M_6$, $m_4(\hat n)$, $m$ & six-dimensional mass, mass read by a clock, Dirac mass (an input) & 3, 10\\
$P_*$, $V_*$; $P_{\rm fam}$ & chiral projector and matter block; family projector & 6\\
$\mathcal M$, $\widetilde{\mathcal C}$, $\Phi$ & kernel space of mass kernels, its complex structure, the kernel field & 7\\
$\Pi$, $\Pi_{\pm A}$, $\Pi_{\pm W}$; $B_{\rm em}$, $Q_A$ & vacuum projectors and aligned points; electric and normal charges & 8\\
$\cR$, $\mathcal G$ & curvature operator on the frame, Einstein tensor & 11\\
\bottomrule
\end{tabular}
\caption{Symbols with one meaning throughout the manuscript.}
\label{tab:fm:notation}
\end{table}

\subsection{Organisation and dependencies}
\label{sec:fm:organisation}

\Cref{ch:ch01_clifford,ch:ch02_lorentz,ch:ch03_readings} develop one constituent: the algebra, the Lorentz structures and complex units with
the clock reduction, and the two readings of $\R^{3,3}$. \Cref{ch:ch04_rigidbody} treats the interior dynamics and its rigid-body
realisation. \Cref{ch:ch05_tower,ch:ch06_chiral,ch:ch07_vacuum,ch:ch08_photon} climb the tower: duplication, the chiral block, the mass kernels
and the vacuum, and the photon. \Cref{ch:ch09_bfield,ch:ch10_classical} treat the dynamics between constituents and the classical and
observational readings. \Cref{ch:ch11_curvature,ch:ch12_causal} pass to curvature and causal structure. \Cref{ch:ch13_compatibility} holds the
sectors against one another, \cref{ch:ch14_related} compares with related approaches, \cref{ch:ch15_methodology} describes the AI-assisted
workflow and the verification, and \cref{ch:ch16_conclusions} concludes. \Cref{app:realisation} gives the explicit matrices, so that every
finite-dimensional claim can be checked by hand; the appendices also hold the dictionaries, the branching tables and the verification
manifest. \Cref{fig:depmap} shows which chapters each chapter relies on.

\begin{figure}[t]
\centering
\begin{tikzpicture}[font=\footnotesize,>=Stealth,
  ch/.style={draw,rounded corners=2pt,align=center,text width=1.85cm,minimum height=0.8cm,inner sep=2pt,fill=white},
  dep/.style={->,thick,mgray}]
\node[ch] (c11) at (5.2,1.5) {11 curvature};
\node[ch] (c12) at (7.8,1.5) {12 causal reduction};
\node[ch] (c14) at (13,1.5) {14 related work};
\node[ch] (c1) at (0,0) {1 Clifford geometry};
\node[ch] (c2) at (2.6,0) {2 Lorentz, clocks};
\node[ch] (c3) at (5.2,0) {3 two readings};
\node[ch] (c4) at (7.8,0) {4 interior, rigid body};
\node[ch] (c10) at (10.4,0) {10 classical};
\node[ch] (c13) at (13,0) {13 compatibility};
\node[ch] (c5) at (0,-1.5) {5 tower};
\node[ch] (c6) at (2.6,-1.5) {6 chiral matter};
\node[ch] (c7) at (5.2,-1.5) {7 kernels, vacuum};
\node[ch] (c8) at (7.8,-1.5) {8 photon};
\node[ch] (c9) at (10.4,-1.5) {9 $b$-field};
\node[ch] (c16) at (13,-1.5) {16 conclusions};
\node[ch] (c15) at (13,-3.0) {15 methods};
\draw[dep] (c1) -- (c2); \draw[dep] (c2) -- (c3); \draw[dep] (c3) -- (c4); \draw[dep] (c4) -- (c10); \draw[dep] (c10) -- (c13);
\draw[dep] (c1) -- (c5); \draw[dep] (c5) -- (c6); \draw[dep] (c6) -- (c7); \draw[dep] (c7) -- (c8);
\draw[dep] (c3) -- (c8); \draw[dep] (c4) -- (c9);
\draw[dep] (c5.south) to[out=-25,in=-155] (c9.south);
\draw[dep] (c3) -- (c11); \draw[dep] (c11) -- (c12); \draw[dep] (c12) -- (c13);
\draw[dep] (c8) -- (c13); \draw[dep] (c9) -- (c13);
\draw[dep] (c13) -- (c14); \draw[dep] (c13) -- (c16); \draw[dep] (c15) -- (c16);
\end{tikzpicture}
\caption{Dependencies between chapters: an arrow from one chapter to another means that the second uses results of the first. Arrows implied by
two others are omitted, so Chapter 11, for instance, also uses Chapter 2. Chapter 13 draws on all earlier chapters, and Chapter 15 describes the
evidence for every chapter.}
\label{fig:depmap}
\end{figure}

% --- source: rest_energy/s0_intro#0 ((head))


% --- source: rest_energy/s0_intro#1 (The question)
\section{Historical motivation: the rest-energy hypothesis}
\label{s0:sec:question}

Part of the programme began with a proposed reading of the rest-energy relation, quoted below. What survived the tests is stated first; the
assessment is in \cref{sec:resynthesis,sec:dupsynthesis}. In the field-strength reading the undivided energy of a field point is the
duality-invariant $E_6$; its rest part $E_0=\tfrac12|q|$ is the infimum of $E_6$ over boosts, and $E_6^2=E_0^2+|\mathbf S|^2$ is an inherited
identity of a field point, not a particle mass. Schur's lemma fixes the isotropic form of the conversion between the two triples, not the
value of $c$, and $c$ is a conversion, not a symmetry. A mass breaks duality and chirality but no spacetime symmetry. The $b$-field splits or
destabilises the binding but never creates a rate or a mass. None of this derives a mass, and the dictionary below consists of readings under
test, not results.

% --- source: rest_energy/s0_intro#2 (The hypothesis in its author's words)
\subsection{The hypothesis in its author's words}

One strand of the present research programme begins with a proposed interpretation of the rest-energy relation:
\begin{quote}
\emph{The mass--energy relation $E=mc^2$ is to be regarded as a unified expression comprising two complementary components: $c^2$,
representing the symmetric component, and $m$, representing the symmetry-breaking component, hypothesized to arise through the
$b$-particle associated with the Yang--Mills $b$-field.}
\end{quote}
Three clarifications fix the intended reading. First, the energy is to be taken as \emph{one undivided entity}, and the mass is to
emerge through an \emph{internal differentiation} of that unity, before any external interaction is applied. Second, $c^2$ corresponds to
the variable $X$ of the interior dynamics of \cref{ch:ch04_rigidbody}, where the shift operators are $\alpha_\pm=\partial_\tau\pm\kappa G$ with $\kappa=c/\ell$. Third, the energy
is six-dimensional. The $b$-field is the non-abelian gauge potential of Yang and Mills \cite{YangMills54}, whose quanta carry isotopic
spin; Yang and Mills left open whether those quanta are massive.

Several points are settled at once. If $E$ and $c^2$ are invariant under a group then so is $m=E/c^2$, so a mass cannot \emph{be} the
breaking of a symmetry that $E$ and $c^2$ share. Positive four-dimensional mass cannot be obtained from momentum along the surplus times,
since $m_4^2=\mu_6^2-(p_u^2+p_v^2)/c^2$ (\cref{r2:prop:rotate}), and non-orthogonality of $\bE$ and $\bB$ generates no mass (\cref{sec:orth}). The
assessment follows the order the hypothesis asks for: the undivided energy, its differentiation, the dynamics of the differentiation, and the
$b$-field last.

% --- source: rest_energy/s0_intro#3 (The second question: duplication, and the bb-field between m)
\subsection{The second question: duplication, and the \texorpdfstring{$b$}{b}-field between manifolds}

After \cref{sec:energy,sec:massbreaks,sec:restclosure,sec:bfield} had established that, on one manifold, the $b$-field can split or destabilise the binding but never
generate it, a further proposal of the author defines the second half of this strand: \emph{the $b$-particle, induced by the $b$-field
through the interaction of two or more manifolds, is proposed as the mechanism enabling the energy cascade across successive
Clifford-algebraic structures, from $\Cl(3,3)$ to $\Cl(6,6)$ and beyond}. This is a hypothesis, not a result. We read it as two questions:
whether a $b$-field acting \emph{between} copies of the arena can do what it cannot do on one copy, and which constructions and processes
take the undivided energy of one copy to two, and what the duplicate can do that a single copy cannot. \Cref{d5:sec:main}, \Cref{d6:sec:main}, \cref{d7:sec:main} answer them and
\cref{s6:sec:synthesis} collects the answers.

% --- source: rest_energy/s0_intro#4 (The dictionary under test)
\subsection{The dictionary under test}

{\sloppy All constructions are made inside the arena of \cref{ch:ch01_clifford}: $V=\R^{3,3}$ with $\eta=\operatorname{diag}(1,1,1,-1,-1,-1)$, the realisation
of $\Cl(3,3)$ on $\R^8$, the binding root $G$ ($Ge_i=e_{i+3}$, $Ge_{i+3}=-e_i$), which squares to $-\id$, is self-adjoint for $\eta$ and
acts as the Hodge star in the field-strength reading $X=\sum_i(B_ie_i+E_ie_{i+3})$, $z=\bB+i\bE$, $q(z)=z\cdot z$
(\cref{thm:emfield}). The working dictionary, whose entries are conjectural readings, is\par}
\begin{center}
\small
\begin{tabular}{@{}>{\raggedright\arraybackslash}p{3.4cm}>{\raggedright\arraybackslash}p{11.0cm}@{}}
\toprule
hypothesis & candidate object in the arena\\
\midrule
$E$, undivided, six-dimensional & the positive, duality-invariant energy $E_6=\tfrac12|z|^2=\tfrac12(|\bB|^2+|\bE|^2)$ of a frame vector; in
the closure of the force law, the Casimir $|p|^2$\\
$c^2$, the symmetric part & the isotropic form $C=c\,\id$ of the conversion between the temporal and the spatial triple, fixed by Schur's
lemma, which does not fix the value of $c$ (\cref{thm:onec}), and the $\eta$-self-adjoint, factorising part $\kappa G$ of the shift, $(\kappa G)^2=-c^2/\ell^2$\\
$m$, the differentiation & the Lorentz-class invariant $E_0=\tfrac12|q|$, the rest-frame energy density of the field; for Dirac matter the mass
form $T=e_1e_2e_3$\\
$b$, the Yang--Mills field & the $\eta$-antisymmetric (transport) part of the shift, $\End(V)=\so{3,3}\oplus\{X=X^\dagger\}$
(\cref{thm:EAdichotomy}); at three constituents the binding-compatible algebra $\so4_\Delta\cong\su2\oplus\su2$ (\cref{thm:PSbinding})\\
\bottomrule
\end{tabular}
\end{center}
Every entry is an interpretive reading to be tested, not a result; \cref{tab:resynthesis,tab:dupsynthesis} give the verdicts. The claims of
this strand are classified as in \cref{sec:deflation}. The relation $E=mc^2$ itself is inherited; what can be new is only how the arena
realises its three ingredients, and which of the realisations are forced.

\section{Use of AI assistance}
\label{sec:fm:disclosure}

This work was developed by the author over roughly three years in an iterative process with AI systems. The questions and the
programme are the author's. Claude models assisted extensively with candidate theorems, exploratory proofs, alternative constructions,
code, calculations and drafting; ChatGPT assisted with discussion, methodology and referee-style criticism. The author examined,
objected to and corrected the results throughout. No particular model version was in use over the whole period. AI-assisted review is
not peer review, and agreement between models is not proof: each statement of this manuscript carries its own evidence, a proof, an exact
computation, a numerical check with its tolerance, or a citation, as described in \cref{ch:ch15_methodology} and listed in
\cref{app:manifest}. The author is responsible for the manuscript and its conclusions.

% restore chapter-based section numbering after the front matter
\renewcommand{\thesection}{\thechapter.\arabic{section}}
